\documentclass{beamer}

% Theme choice
\usetheme{Madrid} % You can change to e.g., Warsaw, Berlin, CambridgeUS, etc.

% Encoding and font
\usepackage[utf8]{inputenc}
\usepackage[T1]{fontenc}

% Graphics and tables
\usepackage{graphicx}
\usepackage{booktabs}

% Code listings
\usepackage{listings}
\lstset{
    basicstyle=\ttfamily\small,
    keywordstyle=\color{blue},
    commentstyle=\color{gray},
    stringstyle=\color{red},
    breaklines=true,
    frame=single
}

% Math packages
\usepackage{amsmath}
\usepackage{amssymb}

% Colors
\usepackage{xcolor}

% TikZ and PGFPlots
\usepackage{tikz}
\usepackage{pgfplots}
\pgfplotsset{compat=1.18}
\usetikzlibrary{positioning}

% Hyperlinks
\usepackage{hyperref}

% Title information
\title{Week 7: Testing and Quality Assurance}
\author{Your Name}
\institute{Your Institution}
\date{\today}

\begin{document}

\frame{\titlepage}

\begin{frame}[fragile]
    \frametitle{Introduction to Testing and Quality Assurance - Overview}
    \begin{block}{Definition}
        Testing is the process of evaluating software to discover any gaps, errors, or missing requirements relative to its intended purpose.
    \end{block}

    \begin{block}{Significance of Testing}
        \begin{itemize}
            \item Ensures Quality: Identifies defects early in the software development life cycle (SDLC).
            \item Enhances User Satisfaction: Validates that the software meets user needs and is free of significant issues.
            \item Reduces Costs: Catching bugs during the testing phase costs less than fixing them post-deployment.
        \end{itemize}
    \end{block}
\end{frame}

\begin{frame}[fragile]
    \frametitle{Testing and Quality Assurance - Key Concepts}
    \begin{block}{Quality Assurance (QA)}
        QA is a systematic process that ensures quality in the development and execution of processes and the resulting products.
    \end{block}

    \begin{block}{Quality Control (QC)}
        QC involves the operational techniques and activities used to fulfill requirements for quality, focusing on detecting defects in the finished product.
    \end{block}
\end{frame}

\begin{frame}[fragile]
    \frametitle{Testing Types}
    \begin{enumerate}
        \item \textbf{Unit Testing:} Tests individual components or modules. \\
            Example: Testing a function like checking if \texttt{calculateSum(a, b)} returns \texttt{a + b}.
        \item \textbf{Integration Testing:} Verifies that modules work together. \\
            Example: Testing communication between front-end and back-end in a web application.
        \item \textbf{System Testing:} Evaluates the complete system's functionality. \\
            Example: Testing an e-commerce application's full process from selection to checkout.
        \item \textbf{Acceptance Testing:} Determines if the system meets business requirements. \\
            Example: User testing simulating real-world use before launch.
    \end{enumerate}
\end{frame}

\begin{frame}[fragile]
    \frametitle{Importance of Testing in Software Development - Part 1}
    \begin{block}{Understanding the Importance of Testing}
        Testing is a critical component of the software development lifecycle. It verifies that software products meet their intended requirements and function correctly before release.
    \end{block}
\end{frame}

\begin{frame}[fragile]
    \frametitle{Importance of Testing in Software Development - Part 2}
    \begin{block}{Consequences of Inadequate Testing}
        \begin{itemize}
            \item \textbf{Presence of Bugs:} 
                \begin{itemize}
                    \item Bugs are defects or issues that arise in software. Without proper testing, these bugs can remain undetected, leading to unexpected behavior during production.
                    \item \textit{Example:} A financial application with a calculation bug could lead to incorrect financial reporting, causing significant legal issues and loss of trust.
                \end{itemize}

            \item \textbf{User Dissatisfaction:}
                \begin{itemize}
                    \item If users encounter bugs or system failures, their overall satisfaction diminishes. Poor user experiences can harm a company's reputation and result in lost customers.
                    \item \textit{Example:} A social media platform that frequently crashes may see users migrate to more stable alternatives.
                \end{itemize}

            \item \textbf{High Maintenance Costs:}
                \begin{itemize}
                    \item Fixing issues post-release is typically more expensive than addressing them during the development phase. Bugs found after deployment often require urgent patches, impacting timelines and budgets.
                    \item \textit{Example:} A software update that needs to be rolled out to address significant bugs can lead to costly emergency work and strained resources.
                \end{itemize}
        \end{itemize}
    \end{block}
\end{frame}

\begin{frame}[fragile]
    \frametitle{Importance of Testing in Software Development - Part 3}
    \begin{block}{Key Points to Emphasize}
        \begin{itemize}
            \item Testing saves costs in the long run by identifying issues early.
            \item User satisfaction is directly tied to the quality of software products.
            \item A well-tested application can provide a competitive advantage and build user trust.
        \end{itemize}
    \end{block}
    
    \begin{block}{Conclusion}
        Implementing a thorough testing strategy is essential to ensure software quality and reliability. By recognizing the consequences of inadequate testing, developers and companies can prioritize testing practices, ultimately leading to more successful software products.
    \end{block}

    \begin{block}{Engagement Tip}
        Encourage students to think of a recent app they used that had issues. What could have been done differently in terms of testing to prevent that problem?
    \end{block}
\end{frame}

\begin{frame}[fragile]
    \frametitle{Types of Testing - Introduction}
    In software development, testing is a critical phase that ensures the final product is functional, reliable, and meets user expectations. 
    Understanding various testing types is essential to developing a robust quality assurance methodology.
\end{frame}

\begin{frame}[fragile]
    \frametitle{Types of Testing - Key Concepts}
    \begin{itemize}
        \item Various testing types focus on different aspects of the software lifecycle.
        \item Combining testing methods enhances product quality and user satisfaction.
        \item Each type plays a distinct role in ensuring robustness and functionality.
    \end{itemize}
\end{frame}

\begin{frame}[fragile]
    \frametitle{Unit Testing}
    \begin{block}{Definition}
        Tests individual components (functions/methods) of the software in isolation to confirm they work as intended.
    \end{block}
    \begin{itemize}
        \item \textbf{Purpose}: Identify and fix bugs before they become part of larger modules.
        \item \textbf{Example}: Validating the function that adds two numbers in a calculator application.
        \item \textbf{Key Point}: Increases code reliability and simplifies future changes.
    \end{itemize}
\end{frame}

\begin{frame}[fragile]
    \frametitle{Integration Testing}
    \begin{block}{Definition}
        Ensures that different modules or services work together as expected.
    \end{block}
    \begin{itemize}
        \item \textbf{Purpose}: Detect interface defects between integrated components.
        \item \textbf{Example}: Testing user authentication module with the database module.
        \item \textbf{Key Point}: Identifies issues in workflows across multiple components.
    \end{itemize}
\end{frame}

\begin{frame}[fragile]
    \frametitle{System Testing}
    \begin{block}{Definition}
        Tests the complete integrated software product to evaluate compliance with specified requirements.
    \end{block}
    \begin{itemize}
        \item \textbf{Purpose}: Validate the overall behavior of the application in the intended environment.
        \item \textbf{Example}: Executing use-case scenarios on an e-commerce platform.
        \item \textbf{Key Point}: Encompasses both functional and non-functional testing aspects.
    \end{itemize}
\end{frame}

\begin{frame}[fragile]
    \frametitle{Acceptance Testing}
    \begin{block}{Definition}
        Conducted to determine if the system meets business needs for final approval.
    \end{block}
    \begin{itemize}
        \item \textbf{Purpose}: Validate the end-to-end business flow and ensure software readiness for release.
        \item \textbf{Example}: Clients testing features in a prototype before final approval.
        \item \textbf{Key Point}: Typically involves end-users verifying that the software is ready for production.
    \end{itemize}
\end{frame}

\begin{frame}[fragile]
    \frametitle{Summary and Additional Considerations}
    \begin{itemize}
        \item Document tests, results, and any relevant findings for future improvements.
        \item Incorporate automated testing tools to streamline the process and maximize efficiency.
    \end{itemize}
    By understanding these testing types, developers can implement effective quality assurance strategies leading to better software delivery.
\end{frame}

\begin{frame}[fragile]
    \frametitle{Unit Testing - Definition and Importance}
    \begin{block}{Definition}
        \textbf{Unit Testing} is a software testing technique where individual components or modules 
        of a software application are tested in isolation to ensure that they function correctly.
    \end{block}
    \begin{itemize}
        \item A "unit" can be a single function, method, or class.
    \end{itemize}

    \begin{block}{Importance of Unit Testing}
        \begin{enumerate}
            \item Catches Bugs Early: Identify and fix bugs before they propagate.
            \item Improves Code Quality: Encourages better design and modularization.
            \item Facilitates Changes: Allows for quick verification of changes.
            \item Documentation: Serves as documentation with usage examples.
        \end{enumerate}
    \end{block}
\end{frame}

\begin{frame}[fragile]
    \frametitle{Unit Testing - Best Practices}
    \begin{itemize}
        \item \textbf{Keep Tests Small and Focused:} Cover a specific aspect of functionality.
        \item \textbf{Use Descriptive Names:} Name tests based on their behavior or expected outcome.
        \item \textbf{Automate Tests:} Use frameworks (e.g., JUnit, pytest) for automatic test execution.
        \item \textbf{Regularly Maintain Test Cases:} Update tests to remain relevant and effective.
    \end{itemize}
\end{frame}

\begin{frame}[fragile]
    \frametitle{Unit Testing - Methods and Examples}
    \begin{block}{Test-Driven Development (TDD)}
        \begin{itemize}
            \item Write the test before writing the actual code.
            \item Fosters a better understanding of requirements.
        \end{itemize}
    \end{block}
    
    \begin{lstlisting}[language=Python, caption={Example of Unit Test in Python}]
    def add(a, b):
        return a + b

    def test_add():
        assert add(2, 3) == 5
        assert add(-1, 1) == 0
    \end{lstlisting}

    \begin{block}{Mocking}
        \begin{itemize}
            \item Use mocks/stubs to isolate the unit of work.
        \end{itemize}
    \end{block}
    
    \begin{lstlisting}[language=Python, caption={Example of Mocking in Python}]
    from unittest.mock import MagicMock

    def fetch_data(api_client):
        response = api_client.get_data()
        return response.json()

    def test_fetch_data():
        mock_client = MagicMock()
        mock_client.get_data.return_value = MagicMock(json=lambda: {'key': 'value'})
        result = fetch_data(mock_client)
        assert result == {'key': 'value'}
    \end{lstlisting}
\end{frame}

\begin{frame}[fragile]
    \frametitle{Integration Testing - Overview}
    \begin{block}{What is Integration Testing?}
        Integration testing is a crucial phase in the software testing lifecycle where individual components or systems are combined and tested as a group. Its primary goal is to evaluate how different modules interact with one another after integration, ensuring they work together seamlessly in a working environment.
    \end{block}
    
    \begin{block}{Purpose of Integration Testing}
        \begin{itemize}
            \item \textbf{Identify Interface Issues:} Uncovers defects occurring during module interactions.
            \item \textbf{Validate Data Flow:} Verifies proper data transmission between modules.
            \item \textbf{Ensure System Cohesiveness:} Confirms that integrated components function correctly as a whole.
        \end{itemize}
    \end{block}
\end{frame}

\begin{frame}[fragile]
    \frametitle{Integration Testing - Key Benefits}
    \begin{itemize}
        \item \textbf{Early Detection of Errors:} Reduces later fixing costs by catching issues earlier.
        \item \textbf{Improves System Quality:} Ensures parts of the system work together, enhancing software quality.
        \item \textbf{Supports Continuous Integration:} Essential for CI/CD pipelines, facilitating faster system updates.
    \end{itemize}
\end{frame}

\begin{frame}[fragile]
    \frametitle{Integration Testing - Methodologies}
    \begin{enumerate}
        \item \textbf{Big Bang Integration Testing}
            \begin{itemize}
                \item \textit{Description:} All components are integrated simultaneously and tested together.
                \item \textit{Pros:} Simple implementation; all components tested in one effort.
                \item \textit{Cons:} Difficult to isolate and debug failures.
            \end{itemize}
            
        \item \textbf{Incremental Integration Testing}
            \begin{itemize}
                \item \textit{Description:} Components are integrated and tested sequentially (Top-Down and Bottom-Up approaches).
                \item \textit{Pros:} Easier defect location, early feedback.
                \item \textit{Cons:} Time-consuming, detailed planning required.
            \end{itemize}
            
        \item \textbf{Sandwich/Hybrid Integration Testing}
            \begin{itemize}
                \item \textit{Description:} Combines top-down and bottom-up approaches.
                \item \textit{Pros:} Flexible and quick at uncovering defects at different levels.
                \item \textit{Cons:} Increased complexity in managing both approaches.
            \end{itemize}
    \end{enumerate}
\end{frame}

\begin{frame}
    \frametitle{Testing Frameworks and Tools}
    Automated testing frameworks streamline the testing process for software applications, enabling developers to write tests that automatically verify functionality, ensuring efficiency and reliability in development.
\end{frame}

\begin{frame}
    \frametitle{Popular Testing Frameworks}
    \begin{block}{1. JUnit for Java}
        \begin{itemize}
            \item \textbf{Overview}: A widely used testing framework for unit testing in Java.
            \item \textbf{Features}:
                \begin{itemize}
                    \item \textbf{Annotations}: Simplifies writing tests with `@Test`, `@Before`, `@After`.
                    \item \textbf{Assertions}: Validates expected outcomes using methods like `assertEquals`, `assertTrue`.
                \end{itemize}
        \end{itemize}
    \end{block}
\end{frame}

\begin{frame}[fragile]
    \frametitle{JUnit Example}
    \begin{lstlisting}[language=Java]
import static org.junit.Assert.*;
import org.junit.Test;

public class ExampleTest {
    @Test
    public void testAddition() {
        assertEquals(5, 2 + 3);
    }
}
    \end{lstlisting}
\end{frame}

\begin{frame}
    \frametitle{Popular Testing Frameworks (cont.)}
    \begin{block}{2. Pytest for Python}
        \begin{itemize}
            \item \textbf{Overview}: A flexible framework supporting simple unit and complex functional tests in Python.
            \item \textbf{Features}:
                \begin{itemize}
                    \item \textbf{Simple Syntax}: Tests can be written as simple functions.
                    \item \textbf{Fixtures}: Supports reusable setup code across tests.
                \end{itemize}
        \end{itemize}
    \end{block}
\end{frame}

\begin{frame}[fragile]
    \frametitle{Pytest Example}
    \begin{lstlisting}[language=Python]
import pytest

def test_addition():
    assert 2 + 3 == 5
    \end{lstlisting}
\end{frame}

\begin{frame}
    \frametitle{Advantages of Automation}
    \begin{enumerate}
        \item \textbf{Efficiency}: Automated tests run faster than manual tests.
        \item \textbf{Repeatability}: Tests can be executed as often as needed.
        \item \textbf{Immediate Feedback}: Quick bug identification reduces costs.
        \item \textbf{Regression Prevention}: Ensures new changes do not break functionality.
    \end{enumerate}
\end{frame}

\begin{frame}
    \frametitle{Key Points and Summary}
    \begin{block}{Key Points}
        \begin{itemize}
            \item Automated frameworks like JUnit and pytest are crucial in modern development.
            \item Integrating automated tests improves software quality.
        \end{itemize}
    \end{block}
    Understanding these frameworks enhances code quality and promotes best practices in software development.
\end{frame}

\begin{frame}[fragile]
    \frametitle{Quality Assurance Best Practices - Introduction}
    \begin{block}{What is Quality Assurance?}
        Quality Assurance (QA) is a systematic process that ensures software products meet specified requirements and standards. 
    \end{block}
    \begin{block}{Importance of Best Practices in QA}
        Establishing best practices in QA contributes to the reliability and performance of software.
    \end{block}
\end{frame}

\begin{frame}[fragile]
    \frametitle{Quality Assurance Best Practices - CI/CD}
    \begin{block}{Continuous Integration (CI) and Continuous Deployment (CD)}
        \begin{itemize}
            \item \textbf{Definition}: CI refers to the frequent integration of code changes into a shared repository. CD automates the deployment of these changes to production.
            \item \textbf{Key Benefits}:
            \begin{itemize}
                \item Early Detection of Defects: Bugs are identified and resolved quickly.
                \item Faster Release Cycles: Automates deployment speeds up the release of new features and fixes.
                \item Example: In a CI/CD pipeline, tools like Jenkins or CircleCI run tests every time new code is pushed.
            \end{itemize}
        \end{itemize}
    \end{block}
\end{frame}

\begin{frame}[fragile]
    \frametitle{Quality Assurance Best Practices - Documentation and Standards}
    \begin{block}{Effective Documentation}
        \begin{itemize}
            \item \textbf{Importance}: Maintains a clear understanding of code, testing procedures, and QA processes.
            \item \textbf{Types of Documentation}:
            \begin{itemize}
                \item Requirement Specifications
                \item Test Plans
                \item User Manuals
            \end{itemize}
            \item Best Practice: Regularly update documentation to reflect changes.
        \end{itemize}
    \end{block}
    
    \begin{block}{Coding Standards}
        \begin{itemize}
            \item \textbf{Definition}: Guidelines to improve code quality, maintainability, and readability.
            \item \textbf{Key Focus Areas}:
            \begin{itemize}
                \item Naming Conventions
                \item Code Structuring
                \item Consistent Syntax
            \end{itemize}
            \item Impact: Reduces complexity and facilitates team collaboration.
        \end{itemize}
    \end{block}
\end{frame}

\begin{frame}[fragile]
    \frametitle{Quality Assurance Best Practices - Summary and Conclusion}
    \begin{block}{Summary of Key Points}
        \begin{itemize}
            \item Implement CI/CD for quick feedback and deployment.
            \item Maintain thorough and up-to-date documentation.
            \item Adopt coding standards to enhance code quality.
        \end{itemize}
    \end{block}
    
    \begin{block}{Conclusion}
        Incorporating these best practices in QA fosters a healthier development environment, leading to higher-quality software and more satisfied users.
    \end{block}
\end{frame}

\begin{frame}[fragile]
    \frametitle{Example of a CI/CD Pipeline}
    \begin{lstlisting}
    Code -> CI Server (Build & Test) 
    -> Deployed to Staging 
    -> Automated Tests 
    -> Production Release
    \end{lstlisting}
\end{frame}

\begin{frame}[fragile]
    \frametitle{Common Challenges in Testing - Overview}
    Here are some challenges typically faced during testing:
    \begin{itemize}
        \item Time Constraints
        \item Test Coverage
        \item Maintaining Test Cases
    \end{itemize}
\end{frame}

\begin{frame}[fragile]
    \frametitle{Common Challenges in Testing - Time Constraints}
    \begin{block}{Time Constraints}
        \begin{itemize}
            \item **Explanation**: Testing phases are often compressed due to tight project deadlines, increasing the risk of undetected defects.
            \item **Example**: When a software product must be launched before a major event, testers may only have a few days for comprehensive testing, leading to critical errors going unnoticed.
        \end{itemize}
    \end{block}
\end{frame}

\begin{frame}[fragile]
    \frametitle{Common Challenges in Testing - Test Coverage}
    \begin{block}{Test Coverage}
        \begin{itemize}
            \item **Explanation**: Adequate test coverage ensures all parts of the application are thoroughly tested. Lack of coverage can leave significant portions unverified.
            \item **Example**: If a web application has multiple user paths but testing only covers the most common scenarios, less frequently used paths may fail in production.
            \item **Key Point**: Aim for a balance between test coverage and project timelines; consider risk-based testing.
        \end{itemize}
    \end{block}
\end{frame}

\begin{frame}[fragile]
    \frametitle{Common Challenges in Testing - Maintaining Test Cases}
    \begin{block}{Maintaining Test Cases}
        \begin{itemize}
            \item **Explanation**: Outdated test cases due to software evolution can lead to false confidence in quality assurance.
            \item **Example**: Modifications to a feature without updating the associated test case can introduce unnoticed bugs.
            \item **Key Point**: Establish a regular schedule for reviewing and updating test cases to keep them relevant.
        \end{itemize}
    \end{block}
\end{frame}

\begin{frame}[fragile]
    \frametitle{Key Takeaways and Conclusion}
    \begin{itemize}
        \item Prioritize testing vital areas based on risk assessment.
        \item Leverage automated testing tools for better coverage and maintainability.
        \item Keep comprehensive documentation of test cases and their statuses.
    \end{itemize}
    
    **Conclusion**: By understanding and addressing these common challenges, teams can enhance their testing processes and improve software quality.
\end{frame}

\begin{frame}[fragile]
    \frametitle{Collaboration in Testing - Overview}
    Collaboration is vital in software testing to ensure thorough coverage, enhance the quality of tests, and streamline communication among team members. Effective teamwork leads to identifying bugs earlier in the development process, resulting in improved software quality and fewer issues in production.
\end{frame}

\begin{frame}[fragile]
    \frametitle{Importance of Teamwork in Testing}
    \begin{itemize}
        \item \textbf{Diverse Skill Sets:} Every team member contributes unique strengths—developers, testers, UX designers, and product owners collaboratively address different aspects of the software.
        \item \textbf{Shared Responsibility:} Testing becomes a collective effort, reducing individual pressure and encouraging thorough peer review of test cases.
        \item \textbf{Improved Communication:} Open communication makes it easier to relay issues or observations quickly, ensuring everyone understands testing goals and progress.
    \end{itemize}
\end{frame}

\begin{frame}[fragile]
    \frametitle{Tools for Collaborative Testing - Git for Version Control}
    \begin{block}{What is Git?}
        Git is a distributed version control system that helps teams manage changes to their codebase, especially useful in collaborative testing settings.
    \end{block}
    \begin{itemize}
        \item \textbf{Key Features:}
        \begin{itemize}
            \item \textbf{Branching:} Work on features or fixes in isolation.
            \item \textbf{Merging:} Combine changes and resolve conflicts.
            \item \textbf{History Tracking:} Full history of changes to revert if needed.
        \end{itemize}
        \item \textbf{Example Workflow:}
        \begin{enumerate}
            \item Create a branch for feature tests (e.g., $feature/user-login$).
            \item Write tests and modify cases in your branch.
            \item Commit changes with meaningful comments.
            \item Push your branch and create a pull request for review.
            \item Conduct code reviews and approve before merging.
        \end{enumerate}
    \end{itemize}
\end{frame}

\begin{frame}[fragile]
    \frametitle{Conclusion - Importance of Rigorous Testing Protocols}
    \begin{enumerate}
        \item \textbf{Quality Assurance}: 
        Rigorous testing protocols ensure software meets specifications and performance standards.
        \begin{itemize}
            \item Example: In a banking app, a bug can prevent transaction completion, eroding user trust.
        \end{itemize}
        
        \item \textbf{Reduction of Costs}:
        Upfront investment in testing reduces long-term maintenance costs.
        \begin{itemize}
            \item Illustration: Fixing a defect during design costs 10 times less than post-deployment fixes.
        \end{itemize}
    \end{enumerate}
\end{frame}

\begin{frame}[fragile]
    \frametitle{Conclusion - Importance of Rigorous Testing Protocols (cont.)}
    \begin{enumerate}
        \setcounter{enumi}{2} % Continue from previous frame
        \item \textbf{User Satisfaction}:
        High-quality software results in greater user satisfaction and retention.
        \begin{itemize}
            \item Example: A well-tested mobile app generates positive reviews and higher engagement.
        \end{itemize}

        \item \textbf{Compliance and Security}:
        Ensures compliance with regulations and enhances security.
        \begin{itemize}
            \item Illustration: Healthcare applications must be tested for HIPAA compliance to protect patient data.
        \end{itemize}

        \item \textbf{Continuous Improvement}:
        Promotes a culture of learning and better practices within development teams.
    \end{enumerate}
\end{frame}

\begin{frame}[fragile]
    \frametitle{Conclusion - Key Points and Summary}
    \begin{block}{Key Points to Emphasize}
        \begin{itemize}
            \item \textbf{Investing in Quality}: Rigorous testing is a strategic investment.
            \item \textbf{Cross-functional Collaboration}: Essential for comprehensive testing.
            \item \textbf{Iterative Testing}: Important throughout the software lifecycle.
        \end{itemize}
    \end{block}

    In summary, rigorous testing protocols are essential for high-quality, reliable, secure, and user-friendly software. Investing in testing reduces costs and enhances customer satisfaction, enabling continuous improvement. Encourage students to reflect on integrating these practices in their future careers.
\end{frame}

\begin{frame}[fragile]
    \frametitle{Conclusion - Call to Action}
    Encourage students to:
    \begin{itemize}
        \item Reflect on integrating rigorous testing protocols into their development practices.
        \item Consider the long-term benefits of quality assurance in their future careers.
    \end{itemize}
\end{frame}


\end{document}